\section{Scientific case studies: the far end of the axis}
\label{sec:cases}

Benchmarks measure what they were built to measure. To see the behaviour above
outside benchmark conditions, we ran AutoDS-Tools on three published scientific
studies, each with an author-reported baseline on the authors' own split.

These runs are the opposite of a knowledge-free condition, which matters for
reading them. Their instructions sit at the far end of the axis of
Section~\ref{sec:axis}: each names the library under a \textsc{required} heading,
supplies a working snippet that constructs it, identifies the target-leaking
columns, describes the split, and states the author baseline as the number to
beat; the polymer instruction goes further and restates the published finding
itself. What they measure is the ceiling under \emph{maximal and correct}
prescription---useful precisely because Section~\ref{sec:mechanism} shows
prescription is not always correct.

\paragraph{Maize yield prediction.} Grain yield for field plots from the
Genomes-to-Fields initiative~\citep{kick2023yield}, on the authors' own split
with no environment leakage. The agent reached Pearson $r = 0.569$ against the authors' best model
at $r = 0.461$ and their neural network at $r = 0.426$. The correlation is
higher; the error is not. The agent's normalised RMSE is $0.958$
against the authors' $0.948$, and its $R^2$ on the held-out set is $0.083$. It
therefore ranks the plots better than the published model while predicting the
yield itself no more accurately, a distinction that matters for the applied
question and that a single-metric comparison would hide.

\paragraph{F-DATA exit-state prediction.} Job success or failure on Fugaku
telemetry~\citep{antici2025fdata}, using only attributes known at submission
time, with the temporal split taken from the paper. The agent reached $92.07\%$ accuracy against an
author baseline of $89\%$. The instruction states the relevant context and the
verifier confirms it: the classes are imbalanced and a constant
``always success'' predictor already scores $88.08\%$. Measured against that
floor rather than against zero, the authors gain $0.9$ points and the agent
$4.0$. Balanced accuracy is $0.705$, so most of the remaining error is on the
minority class, which is the class an operator would want predicted.

\paragraph{OpenPoly glass transition temperature.} Polymer property regression
from PSMILES strings~\citep{wang2025openpoly}. The agent reached $R^2 = 0.543$ against an author baseline
of $0.904$, with MAE $40.4$\,K against $14.6$\,K. This comparison is \emph{not}
like-for-like: the authors used the full dataset, we were unable to obtain it,
and the agent worked from a subset whose test split is $89$ rows. We report the
row for completeness and draw no conclusion from it.

\paragraph{What this does and does not show.} Two of three author baselines were
exceeded on the metric the authors reported, by a $31$B open model. That is a
real result about applied scientific baselines, which are frequently well below
what a competent tabular pipeline achieves, and it supports reading the ceiling
of Section~\ref{sec:ceiling} as a statement about the gap to \emph{tuned}
tabular baselines specifically rather than about the systems being weak in
general. It is not a result about agents working unaided: the knowledge that
would otherwise have to be discovered, which library, which featurisation and
which columns leak, was supplied in the prompt. Read against
Section~\ref{sec:mechanism} the two findings agree. Prescription pays where the
model's own default would be wrong, and on an unfamiliar modality such as
polymer SMILES the default would certainly be wrong; the tabular benchmark
tasks, where the default is already right, are where the same prescription buys
nothing.

These are three single points and not a design. Running the other two systems on
the same cases, once with these instructions and once stripped to the task
statement, would make it a two-by-three crossing and would test
Section~\ref{sec:kdiscterm} on the modality where prescription actually pays. We
did not: unlike the benchmark tasks, these exist only as the runs that produced
them, with no task definitions, splits or scorers in the environment we rebuilt.
